\documentclass[11pt,a4paper]{article}
\usepackage{fontspec,amsmath,amssymb,amsthm,microtype}
\setmainfont{lmroman10-regular.otf}[BoldFont=lmroman10-bold.otf,ItalicFont=lmroman10-italic.otf,BoldItalicFont=lmroman10-bolditalic.otf]
\usepackage{xcolor}
\usepackage[most]{tcolorbox}
\usepackage{mathrsfs}
\providecommand{\lt}{<}
\providecommand{\gt}{>}
\setmonofont{DejaVuSansMono.ttf}[BoldFont=DejaVuSansMono-Bold.ttf,ItalicFont=DejaVuSansMono-Oblique.ttf,BoldItalicFont=DejaVuSansMono-BoldOblique.ttf,Scale=0.86]
\newfontfamily\leanfallback{STIXTwoMath-Regular.otf}[Scale=0.86]
\definecolor{leanrule}{gray}{0.62}
\definecolor{leanback}{gray}{0.962}
\definecolor{leantitle}{gray}{0.22}
\definecolor{leanmuted}{gray}{0.38}
\newtcolorbox{leanbox}[1]{enhanced,breakable,colback=leanback,colframe=leanrule,
  boxrule=0.4pt,arc=0.6mm,left=2.4mm,right=2.4mm,top=1.8mm,bottom=1.8mm,
  before skip=0.8em,after skip=1em,fonttitle=\small\bfseries,coltitle=leantitle,
  colbacktitle=leanback,titlerule=0.3pt,title={#1}}
\newlength{\leancharwidth}
\newenvironment{leancode}{\par\smallskip\begingroup\ttfamily\footnotesize\raggedright
  \settowidth{\leancharwidth}{M}\setlength{\parskip}{0pt}\setlength{\parindent}{0pt}
  \frenchspacing}{\par\endgroup\smallskip}
\newcommand{\leanline}[2]{\hangindent=\dimexpr#1\leancharwidth+4\leancharwidth\relax
  \hangafter=1\noindent\hspace*{#1\leancharwidth}#2\par}
\newcommand{\leansource}[1]{\par{\footnotesize\color{leanmuted}#1}\par}
\newcommand{\leanlabel}[1]{\par\smallskip\noindent{\small\bfseries #1}\par}
\usepackage[margin=25mm]{geometry}
\usepackage[unicode,hidelinks]{hyperref}
\hypersetup{pdfauthor={LangevinLean Project}}
\setlength{\emergencystretch}{3em}
\setlength{\parskip}{0.4em}
\setlength{\parindent}{0pt}
\allowdisplaybreaks[2]
\raggedbottom
\title{\texorpdfstring{The quadratic upper bound}{The quadratic upper bound}}
\author{LangevinLean Project}
\date{}
\begin{document}\maketitle
Draft mathematical exposition; these notes may contain mistakes.


Let \(H\) be a real Hilbert space, let \(L\ge0\), and let \(f:H\to\mathbb R\) be differentiable at every point, with gradient \(\nabla f\) (the vector representing the derivative, \(Df(x)[u]=\langle\nabla f(x),u\rangle\)). The gradient is \emph{\(L\)-Lipschitz} if \(\|\nabla f(u)-\nabla f(v)\|\le L\|u-v\|\) for all \(u,v\in H\). This section proves the quadratic upper bound


\[f(y)\le f(x)+\langle\nabla f(x),y-x\rangle+\frac L2\|y-x\|^2\qquad(x,y\in H),\]


often called the descent lemma. Together with the quadratic lower bound of strong convexity it yields the interpolation inequality behind the sharp contraction of the gradient step, and it is also used for the moment bounds of the Gibbs measure. The proof restricts \(f\) to the segment from \(x\) to \(y\), subtracts the linear part and the quadratic correction, and shows that the resulting function of one variable is nonincreasing; no second derivative of \(f\) is used.


\section{\texorpdfstring{The quadratic upper bound for a Lipschitz gradient}{The quadratic upper bound for a Lipschitz gradient}}


\textbf{Claim.} Let \(H\) be a real Hilbert space and \(L\ge0\). Let \(f:H\to\mathbb R\) be differentiable at every point and assume \(\|\nabla f(u)-\nabla f(v)\|\le L\|u-v\|\) for all \(u,v\in H\). Then for all \(x,y\in H\),


\[f(y)\le f(x)+\langle\nabla f(x),y-x\rangle+\frac L2\|y-x\|^2.\]


\textbf{Proof.} Fix \(x,y\) and put \(d=y-x\), \(a=\langle\nabla f(x),d\rangle\), \(b=\|d\|^2\). Define \(q:\mathbb R\to\mathbb R\) by


\[q(t)=f(x+td)-f(x)-ta-\frac L2t^2b.\]


Then \(q(0)=0\) and, since \(x+d=y\), \(q(1)=f(y)-f(x)-\langle\nabla f(x),y-x\rangle-\frac L2\|y-x\|^2\). The claim is therefore the inequality \(q(1)\le q(0)\).




\emph{(a) Derivative of \(q\).} The curve \(t\mapsto x+td\) has derivative \(d\) at every \(t\). By the chain rule and the definition of the gradient,


\[\frac{\mathrm d}{\mathrm dt}f(x+td)=Df(x+td)[d]=\langle\nabla f(x+td),d\rangle\qquad(t\in\mathbb R).\]


The remaining terms of \(q\) are polynomials in \(t\), with derivatives \(a\) and \(\frac L2\cdot2t\,b=Ltb\). Hence \(q\) is differentiable on \(\mathbb R\) and


\[q'(t)=\langle\nabla f(x+td),d\rangle-a-Ltb.\]


\emph{(b) \(q'(t)\le0\) for \(0<t<1\).} Let \(0<t<1\); in particular \(t\ge0\), so \(\|td\|=t\|d\|\). The Lipschitz hypothesis at the points \(x+td\) and \(x\) gives


\[\|\nabla f(x+td)-\nabla f(x)\|\le L\|td\|=Lt\|d\|.\qquad(1)\]


By the Cauchy–Schwarz inequality \(\langle u,v\rangle\le\|u\|\|v\|\), and then by (1) multiplied by \(\|d\|\ge0\),


\[\langle\nabla f(x+td),d\rangle-a=\langle\nabla f(x+td)-\nabla f(x),d\rangle\le\|\nabla f(x+td)-\nabla f(x)\|\,\|d\|\le Lt\|d\|^2=Ltb.\]


By (a) this says exactly \(q'(t)\le0\).




\emph{(c) \(q\) is nonincreasing on \([0,1]\).} Being differentiable on \(\mathbb R\), \(q\) is continuous on \([0,1]\) and differentiable on \((0,1)\). Let \(0\le s<t\le1\). By the mean value theorem there is \(\xi\in(s,t)\subseteq(0,1)\) with \(q(t)-q(s)=q'(\xi)(t-s)\), and this is \(\le0\) by (b). Hence \(q(t)\le q(s)\).




\emph{(d) Conclusion.} By (c) with \(s=0\), \(t=1\), we get \(q(1)\le q(0)=0\), that is,


\[f(x+d)-f(x)-a-\frac L2b\le0.\]


Substituting \(x+d=y\), \(a=\langle\nabla f(x),y-x\rangle\) and \(b=\|y-x\|^2\) and rearranging gives the claim. \(\square\)




Equivalently, \(f(y)-f(x)-\langle\nabla f(x),d\rangle=\int_0^1\langle\nabla f(x+td)-\nabla f(x),d\rangle\,dt\le\int_0^1Lt\|d\|^2\,dt=\frac L2\|d\|^2\); the monotonicity argument above avoids the fundamental theorem of calculus for the derivative \(t\mapsto\langle\nabla f(x+td),d\rangle\), whose integrability would otherwise have to be checked.

\begin{leanbox}{Lean: \texttt{smooth\_\allowbreak{}upper\_\allowbreak{}bound}}
\leansource{\texttt{smooth\_\allowbreak{}upper\_\allowbreak{}bound} \quad UnadjustedLangevin/\allowbreak{}Analysis/\allowbreak{}Smooth.lean, lines 27–73}
\begin{leancode}
\leanline{0}{theorem\ smooth\_upper\_bound\ (f\ :\ H\ →\ ℝ)\ (L\ :\ ℝ≥0)}
\leanline{4}{(hf\ :\ Differentiable\ ℝ\ f)\ (hLip\ :\ LipschitzWith\ L\ (gradient\ f))\ (x\ y\ :\ H)\ :}
\leanline{4}{f\ y\ ≤\ f\ x\ +\ inner\ ℝ\ (gradient\ f\ x)\ (y\ -\ x)\ +\ (L\ :\ ℝ)\ /\ 2\ *\ ‖y\ -\ x‖\ \textasciicircum{}\ 2}
\end{leancode}
\leanlabel{How the Lean proof follows the argument}
\texttt{smooth\_\allowbreak{}upper\_\allowbreak{}bound} takes \texttt{L : \ensuremath{\mathbb R}\ensuremath{\ge}0}, a nonnegative real number (the coercion \texttt{(L : \ensuremath{\mathbb R})} is the same number viewed in \(\mathbb R\)), \texttt{hf : Differentiable \ensuremath{\mathbb R} f}, and \texttt{h\allowbreak{}Lip : Lipschitz\allowbreak{}With L (gradient f)}, which is Mathlib's statement \(\operatorname{dist}(\nabla f(u),\nabla f(v))\le L\operatorname{dist}(u,v)\), i.e. the Lipschitz bound of the claim. The space carries \texttt{[Complete\allowbreak{}Space H]} so that \texttt{gradient f} is defined. The proof is the argument above.



\begin{enumerate}


\item The local definitions \texttt{d}, \texttt{a}, \texttt{b} and \texttt{q} are \(d\), \(a\), \(b\) and \(q\).

\item \texttt{hline} is the derivative \(d\) of \(t\mapsto x+td\), obtained from \texttt{has\allowbreak{}Deriv\allowbreak{}At\_\allowbreak{}id} by \texttt{smul\_\allowbreak{}const d} and \texttt{const\_\allowbreak{}add x}. \texttt{hcomp} is the chain rule \texttt{(hf (x + t \ensuremath{\bullet} d)).\allowbreak{}has\allowbreak{}FDeriv\allowbreak{}At.\allowbreak{}comp\_\allowbreak{}has\allowbreak{}Deriv\allowbreak{}At t (hline t)}; the \texttt{simpa} step uses \texttt{inner\_\allowbreak{}gradient\_\allowbreak{}left} to write the derivative as \(\langle\nabla f(x+td),d\rangle\).

\item \texttt{hq} is (a). The term after \texttt{convert} builds the derivative of \(q\) from the rules for subtracting a constant (\texttt{sub\_\allowbreak{}const}), for the linear term (\texttt{(has\allowbreak{}Deriv\allowbreak{}At\_\allowbreak{}id t).\allowbreak{}mul\_\allowbreak{}const a}) and for the quadratic term (\texttt{((has\allowbreak{}Deriv\allowbreak{}At\_\allowbreak{}id t).\allowbreak{}pow 2).\allowbreak{}const\_\allowbreak{}mul} and \texttt{mul\_\allowbreak{}const b}); the two side goals of \texttt{convert .\allowbreak{}.\allowbreak{}.\allowbreak{} using 1} identify the function with \(q\) (\texttt{ext s}, \texttt{simp only}) and the derivative with \(\langle\nabla f(x+td),d\rangle-a-Ltb\) (\texttt{ring}).

\item \texttt{hq\_\allowbreak{}nonpos} is (b) for \(t\) in \texttt{interior (Set.\allowbreak{}Icc (0 : \ensuremath{\mathbb R}) 1)}, the open interval \((0,1)\): \texttt{ht0} gives \(t\ge0\); \texttt{hnorm} is (1), from \texttt{h\allowbreak{}Lip.\allowbreak{}dist\_\allowbreak{}le\_\allowbreak{}mul (x + t \ensuremath{\bullet} d) x} with distances written as norms; \texttt{hinner} is the Cauchy–Schwarz inequality \texttt{real\_\allowbreak{}inner\_\allowbreak{}le\_\allowbreak{}norm}; \texttt{hprod} multiplies \texttt{hnorm} by \(\|d\|\) (\texttt{mul\_\allowbreak{}le\_\allowbreak{}mul\_\allowbreak{}of\_\allowbreak{}nonneg\_\allowbreak{}right}); after \texttt{inner\_\allowbreak{}sub\_\allowbreak{}left}, \texttt{nlinarith} combines them.

\item \texttt{hanti} is (c): the Mathlib lemma \texttt{antitone\allowbreak{}On\_\allowbreak{}of\_\allowbreak{}deriv\_\allowbreak{}nonpos} (a consequence of the mean value theorem) applied on the convex set \texttt{Set.\allowbreak{}Icc 0 1}, with continuity and differentiability obtained from \texttt{hq}, and with \texttt{deriv q t} identified with the formula of \texttt{hq} by \texttt{(hq t).\allowbreak{}deriv}.

\item \texttt{hfinal} is \(q(1)\le q(0)\). The \texttt{dsimp} and \texttt{simp only} calls evaluate \(q\) at \(0\) and \(1\), \texttt{hxy : x + (y -\allowbreak{} x) = y} (proved by \texttt{abel}) replaces \(x+d\) by \(y\), and \texttt{linarith} rearranges as in (d).

\end{enumerate}
\end{leanbox}

\end{document}
